\begin{tikzpicture}[x=1mm,y=1mm,font=\figurefont\footnotesize,
  every node/.style={text=NNInk},
  nn flow/.append style={draw=NNWire,rounded corners=1.5mm},
  nn operator path/.append style={draw=NNWire,rounded corners=1.5mm},
  op/.style={nn operator,minimum width=33mm,inner sep=2pt}]
  \node[op,draw=NNInput,fill=NNInput!18] (in) at (55,0) {“追”的表示 $\bm h_2$};
  \node[op] (value) at (25,-22) {候选投影 $\bm W_b$};
  \node[op] (gate) at (85,-22) {门投影 $\bm W_g$ 与 $\psi$};
  \node[align=center] (v) at (25,-41) {$\bm B_{2,:}$\\$(2,\ {-4})$};
  \node[align=center] (g) at (85,-41) {$\psi(\bm D_{2,:})$\\$(0.5,\ 2)$};
  \node[nn pointwise] (mul) at (55,-58) {$\odot$};
  \node[anchor=west] at (62,-67) {$(1,\ {-8})$};
  \node[nn transform,minimum width=33mm] (out) at (55,-78) {输出投影 $\bm W_o$};
  \draw[nn flow,-,sharp corners] (in.south)--(55,-10);
  \draw[nn flow] (55,-10)--(25,-10)--(value.north);
  \draw[nn flow] (55,-10)--(85,-10)--(gate.north);
  \fill[NNWire] (55,-10) circle[radius=.45mm];
  \draw[nn operator path] (value.south)--(v.north);
  \draw[nn operator path] (gate.south)--(g.north);
  \draw[nn operator path] (v.south)--(25,-58)--(mul.west);
  \draw[nn operator path] (g.south)--(85,-58)--(mul.east);
  \draw[nn flow,draw=NNOutput] (mul.south)--(out.north);
  \node[nn annotation] at (55,-91) {猫、追、狗分别使用同一套参数};
\end{tikzpicture}
